\section{Disjoint Sparse Table}
\whenuse{Arreglo est\'atico, operaci\'on asociativa y muchas consultas de rango
en $O(1)$, incluso cuando la operaci\'on no es idempotente (suma, producto).}

Consultas est\'aticas $O(1)$ para cualquier operaci\'on asociativa, incluso no
idempotente y no conmutativa. Los rangos son inclusivos.

\complexity{build $O(n\log n)$; query $O(1)$; memoria $O(n\log n)$.}

\begin{lstlisting}[style=cpp]
template<class T, class Op = plus<T>>
struct DisjointSparseTable {
    int n = 0, levels = 0;
    Op op;
    vector<T> values;
    vector<vector<T>> table;

    explicit DisjointSparseTable(const vector<T> &a, Op operation = Op())
        : n(a.size()), op(operation), values(a) {
        while ((1 << levels) < max(1, n)) ++levels;
        table.assign(levels, vector<T>(n));
        for (int k = 0; k < levels; ++k) {
            int block = 1 << (k + 1);
            for (int begin = 0; begin < n; begin += block) {
                int mid = min(begin + (block >> 1), n);
                int end = min(begin + block, n);
                if (begin < mid) {
                    table[k][mid - 1] = values[mid - 1];
                    for (int i = mid - 2; i >= begin; --i)
                        table[k][i] = op(values[i], table[k][i + 1]);
                }
                if (mid < end) {
                    table[k][mid] = values[mid];
                    for (int i = mid + 1; i < end; ++i)
                        table[k][i] = op(table[k][i - 1], values[i]);
                }
            }
        }
    }

    T query(int l, int r) const { // [l, r], requiere indices validos
        if (l == r) return values[l];
        if (l > r) swap(l, r);
        int k = 31 - __builtin_clz((unsigned)(l ^ r));
        return op(table[k][l], table[k][r]);
    }
};
\end{lstlisting}
